 \begin{minipage}{0.40\linewidth}
 \includegraphics[scale=0.3]{oshw-logo.eps}
 \end{minipage}\hfill
 \begin{minipage}{0.55\linewidth}
Dans cette partie, nous allons parler des matériels \textit{open-source}. La conception est effectuée de manière communautaire ou par des sociétés privées. Elles ne sont pas soumises à des brevets ultra-restricitfs mais le plus souvent sous licence GPL (General Public License) plus souples. 
 \end{minipage}
Le matériel \textit{open-source} est très peu developpé contrairement à son homologue logiciel. Il n'y a que peu d'acteurs sur cette scène alors qu'il y en a des dizaines pour l'\textit{open-source} logiciel. Cependant, ce pan de l'\textit{open-source} tend de plus en plus à se développer ces derniers temps. Pour preuve, en 2010 a été publié l'\textbf{Open Source Hardware} (\textbf{OSHW} qui définit ce qu'est un matériel \textit{open-source} et les idées derrière ce mouvement.

Il y a d'autres projets en cours, que ce soit dans le domaine des ordinateurs eux-mêmes, dans leur composants (CPU (\textit{Central Processing Unit}, le coeur de la machine), carte graphique, ...), voire même dans d'autres secteurs tels que la téléphonie et la géolocalisation. Voici la liste des projets \textit{open-source} relatifs aux ordinateurs : \textit{Arduino}, \textit{Beagle Board}, \textit{Bug Labs}, \textit{CuBox}, \textit{Ethernut}, \textit{GP32}, \textit{Maple}, \textit{PLAICE}, \textit{Raspberry Pi}, \textit{Simputer}, \textit{Pandora}, \textit{PC532} \footnote{\url{http://en.wikipedia.org/w/index.php?title=List_of_open-source_hardware_projects\&oldid=533515806}}
qui date des années 1990, \textit{Minimig} qui est une nouvelle implémentation de l'Amiga 500 utilisant des FPGA, \textit{Milkymist One} qui est un SoC évolué implémentant tous les principes logiques à travers des fichiers de conception créés avec \textit{Verilog}. Ces fichiers sont \textit{open-source}, donc disponibles pour tout le monde. Dans la plupart des autres matériels \textit{open-source}, ces fichiers sont confidentiels et propriétaires.

Nous pouvons lister différents projets : \textit{Amber}, \textit{LEON}, \textit{OpenCores}, \textit{OpenSPARC}, \textit{OpenRISC} qui est plus un groupe de développeurs visant à créer des CPU RISC ayant des performances très importantes. Enfin, le \textit{Milkymist SoC}, couplé au micro-processeur \textit{LatticeMico32}, qui possède un ensemble de périphériques développés de manière indépendante et \textit{open-source}.

La plus grande communauté à l'heure actuelle est \textit{OpenCores} (\url{http://www.opencores.org}). Fondée en 1999 par Damjan Lampret, \textit{OpenCores} a developpé \textit{OpenRISC 1000}. Le développement s'est poursuivi jusqu'en 2007. Durant cette période \textit{OpenCores} a été supporté par \textit{Flextronics}. À partir de 2007, \textit{OpenRISC 1000} est passé dans une phase de déploiement commercial grâce à \textit{ORSoC AB}. En mars 2012, il y avait plus de 147 000 utilisateurs enregistrés et plus de 900 projets en cours sur cette architecture matérielle. \textit{ORPSoC} est l'implémentation de référence sur cette architecture. Cependant, le \textit{Raspberry Pi}, l'ordinateur à 29\$, sorti courant 2012, grandit de plus en plus. De nombreuses cartes électroniques complémentaire au \textit{Raspberry Pi} sont sorties, toutes à un prix raisonnable.


Il faut savoir que contrairement aux idées reçues, les entreprises qui font de l'\textit{open-source hardware} ont de bons résultats. En effet, les 13 plus grosses sociétés qui ont ce genre d'activités ont toutes plus d'un million de dollars de chiffre d'affaire (en dollars et en 2010)\cite{TORRONE_FRIED}:
 \begin{itemize}
 \item Adafruit Industries : plus d'1 million
 \item Arduino : plus d'1 million
 \item BeagleBoard : 1 million
 \item Bug Labs : 1 million
 \item Chumby : plus d'1 million
 \item Dangerous Prototypes
 \item DIY Drones : 1 million
 \item Evil Mad Scientist  Labs
 \item Liquidware : plus d'1 million
 \item MakerBot Industries : 1 million
 \item Maker Shed : plus d'1 million
 \item Parallax : plus d'1 million
 \item Seeed Studios : 1 million
 \item Solarbotics : plus d'1 million
 \item SparkFun Electronics : 10 millions
 \end{itemize}  
Ces entreprises ne sont qu'un échantillon. Il en existe bien d'autres de ce type qui approchent ou dépassent le million de chiffre d'affaire. De plus, ces chiffres datent de 2010, on peut donc légitimement supposer que ces chiffres ont évolué de manière positive ces deux dernières années. Selon les prévisions pour 2015, le chiffre d'affaire cumulé de toutes les entreprises de matériel \textit{open-source} devrait avoisiner le milliard de dollars.

Ce secteur a donc un avenir radieux qui s'offre à lui. De plus en plus de projets voient le jour et ceux déjà présents s'améliorent grâce à la contribution de personnes possédant ce genre de matériel car c'est bien là le leitmotiv de l'\textit{open-source} : la progression par la participation.
